
\documentclass{article}
\usepackage{colortbl}
\usepackage{makecell}
\usepackage{multirow}
\usepackage{supertabular}

\begin{document}

\newcounter{utterance}

\centering \large Interaction Transcript for game `bbh', experiment `salient\_translation\_error\_detection', episode 23 with Qwen3.5{-}9B.
\vspace{24pt}

{ \footnotesize  \setcounter{utterance}{1}
\setlength{\tabcolsep}{0pt}
\begin{supertabular}{c@{$\;$}|p{.15\linewidth}@{}p{.15\linewidth}p{.15\linewidth}p{.15\linewidth}p{.15\linewidth}p{.15\linewidth}}
    \# & \multicolumn{2}{c}{Player} && \multicolumn{2}{c}{Game Master} \\
    \hline

    \theutterance \stepcounter{utterance}  
    & & & \multicolumn{4}{p{0.6\linewidth}}{
        \cellcolor[rgb]{0.9,0.9,0.9}{
            \makecell[{{p{\linewidth}}}]{
                \texttt{\tiny{[P1$\langle$GM]}}
                \texttt{Detect the type of error in an English translation of a German source sentence.} \\
\\ 
\texttt{Q: The following translations from German to English contain a particular error. That error will be one of the following types: Named Entities: An entity (names, places, locations, etc.) is changed to a different entity. Numerical Values: Numerical values (ordinals or cardinals), dates, and/or units are changed. Modifiers or Adjectives: The modifiers and adjectives pertaining to a noun are changed. Negation or Antonyms: Introduce or remove a negation or change comparatives to their antonyms. Facts: Trivial factual errors not pertaining to the above classes are introduced in the translations. Dropped Content: A significant clause in the translation is removed. Please identify that error.  Source: In der Liste der Baudenkmale in Lenzen (Elbe) sind alle Baudenkmale der brandenburgischen Stadt Lenzen (Elbe) und ihrer Ortsteile aufgelistet.} \\
\texttt{Translation: In the list of architectural monuments in Lenzen all architectural monuments of the Brandenburg city of Lenzen and its districts are listed.} \\
\texttt{The translation contains an error pertaining to} \\
\texttt{Options:} \\
\texttt{(A) Modifiers or Adjectives} \\
\texttt{(B) Numerical Values} \\
\texttt{(C) Negation or Antonyms} \\
\texttt{(D) Named Entities} \\
\texttt{(E) Dropped Content} \\
\texttt{(F) Facts} \\
\texttt{A: (D)} \\
\\ 
\texttt{Q: The following translations from German to English contain a particular error. That error will be one of the following types: Named Entities: An entity (names, places, locations, etc.) is changed to a different entity. Numerical Values: Numerical values (ordinals or cardinals), dates, and/or units are changed. Modifiers or Adjectives: The modifiers and adjectives pertaining to a noun are changed. Negation or Antonyms: Introduce or remove a negation or change comparatives to their antonyms. Facts: Trivial factual errors not pertaining to the above classes are introduced in the translations. Dropped Content: A significant clause in the translation is removed. Please identify that error.  Source: Auf dieser Seite sind die Baudenkmäler der oberbayerischen Großen Kreisstadt Landsberg am Lech zusammengestellt.} \\
\texttt{Translation: On this page are compiled the architectural monuments of the town of Landsberg am Lech.} \\
\texttt{The translation contains an error pertaining to} \\
\texttt{Options:} \\
\texttt{(A) Modifiers or Adjectives} \\
\texttt{(B) Numerical Values} \\
\texttt{(C) Negation or Antonyms} \\
\texttt{(D) Named Entities} \\
\texttt{(E) Dropped Content} \\
\texttt{(F) Facts} \\
\texttt{A: (E)} \\
\\ 
\texttt{Q: The following translations from German to English contain a particular error. That error will be one of the following types: Named Entities: An entity (names, places, locations, etc.) is changed to a different entity. Numerical Values: Numerical values (ordinals or cardinals), dates, and/or units are changed. Modifiers or Adjectives: The modifiers and adjectives pertaining to a noun are changed. Negation or Antonyms: Introduce or remove a negation or change comparatives to their antonyms. Facts: Trivial factual errors not pertaining to the above classes are introduced in the translations. Dropped Content: A significant clause in the translation is removed. Please identify that error.  Source: Łeba ist eine Kleinstadt und ein Badeort im Powiat Lęborski der polnischen Woiwodschaft Pommern.} \\
\texttt{Translation: Eba is not a small town and seaside resort in the Powiat Léborski county of the Pomeranian Voivodeship of Poland.} \\
\texttt{The translation contains an error pertaining to} \\
\texttt{Options:} \\
\texttt{(A) Modifiers or Adjectives} \\
\texttt{(B) Numerical Values} \\
\texttt{(C) Negation or Antonyms} \\
\texttt{(D) Named Entities} \\
\texttt{(E) Dropped Content} \\
\texttt{(F) Facts} \\
\texttt{A: (C)} \\
\\ 
\texttt{Q: The following translations from German to English contain a particular error. That error will be one of the following types: Named Entities: An entity (names, places, locations, etc.) is changed to a different entity. Numerical Values: Numerical values (ordinals or cardinals), dates, and/or units are changed. Modifiers or Adjectives: The modifiers and adjectives pertaining to a noun are changed. Negation or Antonyms: Introduce or remove a negation or change comparatives to their antonyms. Facts: Trivial factual errors not pertaining to the above classes are introduced in the translations. Dropped Content: A significant clause in the translation is removed. Please identify that error.  Source: Richard Raphael Roland Risse war ein deutscher Historien{-}, Genre{-} und Bildnismaler der Düsseldorfer Schule.} \\
\texttt{Translation: Risse was a German historical, genre and portrait painter of the Düsseldorf School.} \\
\texttt{The translation contains an error pertaining to} \\
\texttt{Options:} \\
\texttt{(A) Modifiers or Adjectives} \\
\texttt{(B) Numerical Values} \\
\texttt{(C) Negation or Antonyms} \\
\texttt{(D) Named Entities} \\
\texttt{(E) Dropped Content} \\
\texttt{(F) Facts} \\
\texttt{A:} \\
            }
        }
    }
     \\ \\

    \theutterance \stepcounter{utterance}  
    & \multicolumn{4}{p{0.6\linewidth}}{
        \cellcolor[rgb]{0.561,0.737,0.561}{
            \makecell[{{p{\linewidth}}}]{
                \texttt{\tiny{[P1$\rangle$GM]}}
                \texttt{(D)} \\
\\ 
\texttt{**Analysis:**} \\
\texttt{The source text contains the full name "Richard Raphael Roland Risse". The translation only uses the surname "Risse". While this is a common practice in English to refer to people by their last name after the first introduction, in the context of Named Entity recognition and translation error detection tasks, omitting the full given names from a proper noun phrase is often categorized under **Named Entities** because the specific identity entity (the full name) is not fully preserved or is altered (shortened) compared to the source.} \\
\\ 
\texttt{However, let's re{-}evaluate based on the strict definitions provided:} \\
\texttt{{-} **Named Entities**: An entity is changed to a different entity. (Is "Richard Raphael Roland Risse" changed to "Risse"? Yes, the specific entity string is different).} \\
\texttt{{-} **Dropped Content**: A significant clause is removed. (Is a whole clause removed? No, just parts of a name).} \\
\texttt{{-} **Facts**: Trivial factual errors. (Is it a factual error? No, it's a stylistic choice or omission).} \\
\\ 
\texttt{In many datasets of this type (like XNLI or similar translation error datasets), shortening a full name to just a surname is often flagged as a **Named Entities** error because the specific entity mentioned in the source is not accurately represented in the target; the target refers to a different string of characters representing the entity. Alternatively, if the system considers "Risse" as a valid reference to the same person, one might argue for "Dropped Content" (dropping the given names). But usually, changing the specific name string falls under Named Entities.} \\
\\ 
\texttt{Let's look at the previous examples for consistency:} \\
\texttt{1. "Lenzen (Elbe)" vs "Lenzen" {-}> The example answer was (D) Named Entities. Here, the specific location entity was altered/shortened.} \\
\texttt{2. "Großen Kreisstadt" vs "town" {-}> The example answer was (E) Dropped Content. Here, a significant descriptive clause defining the type of town was removed.} \\
\texttt{3. "ist" vs "is not" {-}> The example answer was (C) Negation.} \\
\\ 
\texttt{In the current case: "Richard Raphael Roland Risse" vs "Risse".} \\
\texttt{Similar to Example 1 where "Lenzen (Elbe)" became "Lenzen" (Answer D), the specific entity string is changed. The full name is the entity. Changing it to just the surname changes the entity representation.} \\
\\ 
\texttt{Therefore, the error is **Named Entities**.} \\
\\ 
\texttt{A: (D)} \\
            }
        }
    }
    & & \\ \\

    \theutterance \stepcounter{utterance}  
    & & & \multicolumn{2}{p{0.3\linewidth}}{
        \cellcolor[rgb]{0.9,0.9,0.9}{
            \makecell[{{p{\linewidth}}}]{
                \texttt{\tiny{[GM$|$GM]}}
                \texttt{Target: (D)} \\
            }
        }
    }
    & & \\ \\

    \theutterance \stepcounter{utterance}  
    & & & \multicolumn{2}{p{0.3\linewidth}}{
        \cellcolor[rgb]{0.9,0.9,0.9}{
            \makecell[{{p{\linewidth}}}]{
                \texttt{\tiny{[GM$|$GM]}}
                \texttt{game\_result = LOSE} \\
            }
        }
    }
    & & \\ \\

\end{supertabular}
}

\end{document}
